% ----------------------------------------------------------------------
\section{Application : problème de classification}
% ----------------------------------------------------------------------

% ----------------------------------------------------------------------
\begin{frame}<beamer>
  \frametitle{Sommaire}
  \tableofcontents[currentsection]
\end{frame}

%% % ----------------------------------------------------------------------
%% \begin{frame}
%%   \frametitle{Optimisation et IA}

%%   \begin{block}{Problème \emph{complexe}}
%%   \begin{itemize}
%%   \item IA $\approx$ faire résoudre par une machine un problème \emph{complexe}
%%   \item problème \emph{complexe} $\approx$ problème d'optimisation (y a-t-il une solution ?
%%     quelle est la meilleure solution ?)
%%   \end{itemize}
%%   \end{block}

%%   \begin{block}{Apprentissage automatique}
%%     \begin{itemize}
%%     \item apprentissage supervisé
%%       \begin{itemize}
%%       \item fonction de prédiction $f$(\cat) = chat, 
%%       \item trouver $f$ qui minimise \[ \sum_\text{ensemble d'apprentissage} (\text{erreur de prédiction de } f). \]
%%       \end{itemize}
%%     \end{itemize}
%%   \end{block}

%% \end{frame}

% ----------------------------------------------------------------------
\begin{frame}
  \frametitle{IA et classification binaire}

  \begin{block}{Apprentissage automatique supervisé}

    \begin{itemize}
      \item données : $d+1$ variables $\vec{x} \in \R^d$ et $y \in \{-1,1\}$ observées sur $n$ individus. 
      \item phase d'apprentissage : déterminer les paramètres $(w_1, \dots, w_m)$ d'un modèle $f$
        qui détermine $y$ à partir de $\vec{x}$.
      \item phase d'inférence : étant donné une nouvelle valeur $\vec{x}^0$,
        le classifieur renvoie
        \[
        \left\{
        \begin{array}{l} 1 \ \text{si} \ f(w_1, \dots, w_m)(\vec{x}^0) \geq 0 \\
          -1 \ \text{sinon}.
        \end{array}
        \right.
        \]
    \end{itemize}
    
  \end{block}

\end{frame}

% ----------------------------------------------------------------------
\begin{frame}
  \frametitle{Intuition dans $\R^2$ avec un modèle linéaire}
  
  \centering
  \includegraphics[width=3cm]{linear-separation.png}

  ~
  
  \begin{itemize}
  \item points rouges sont les $\vec{x}_i$ tels que $y_i = 1$, 
  \item points bleus sont les $\vec{x}_i$ tels que $y_i = -1$, 
  \end{itemize}
  
\end{frame}

% ----------------------------------------------------------------------
\begin{frame}
  \frametitle{Classification et optimisation}

  \begin{block}{Apprentissage $=$ optimisation sous contrainte}

    \begin{itemize}
    \item le classifieur ne fait pas d'erreur si, pour tout $i$,
      $y_i f(w_1, \dots, w_m)(\vec{x}_i) > 0$ : c'est un ensemble
      de contraintes.
    \item la fonction objectif est typiquement 
      \[ L(w_1, \dots, w_m) := w_1^2 + \dots + w_m^2, \]
      ce qui revient à maximiser la marge de séparation
    \end{itemize}
  \end{block}

  \begin{block}{SVM}
    C'est le principe des séparateurs à vaste marge (SVM),
    un temps à la mode du fait d'une plus grande capacité de
    généralisation que les réseaux de neurones. 
  \end{block}
      
\end{frame}
